\subsection{Resolventes del transporte y conservación de la memoria}
\label{subsec:ii-memoria-resolvente}

La historia prolongada determina las ventanas
\(I_m=\{9m+1,\ldots,9m+9\}\), para \(m\geq0\).
Con la numeración de esta sección, \(I_m=E_m\) durante el primer ciclo.
Sea \(\sigma_m\in\{-1,1\}\) la orientación conservada por esa historia;
para \(0\leq m<12\), coincide con \(\Sigma_{12}(m)\). Se define
\[
 a_m=\sigma_m\sum_{t\in I_m}\mathbf1_{\{2,4,6,8\}}(d_t),
 \qquad |a_m|\leq9.
\]
El registro y el marco móvil se prolongan por
\[
 \mathbf z_{m+1}=\mathbf z_m+a_m\mathbf e_{m\bmod12},\qquad y_m=S^{-m}\mathbf z_m.
\]
El estado inicial \(\mathbf z_0\) conserva la historia previa; vale cero para
un prefijo inicialmente vacío. En este apartado escribimos
\(R=S^{-1}\) para el transporte del marco. La actualización
\eqref{eq:eta} se mantiene en toda ventana. La orientación y los eventos
posteriores se obtienen de la historia prolongada: el retorno del marco
no presupone periodicidad de la secuencia de eventos.

\begin{proposition}[Identidad de truncación con estado terminal]
Para $N\geq1$, sean
\[
Y_N(u)=\sum_{m=0}^N y_mu^m,\qquad
H_{\eta,N}(u)=\sum_{m=0}^{N-1}\eta_mu^m.
\]
Se cumple la identidad polinómica exacta
\begin{equation}
(I-uR)Y_N(u)=y_0+uH_{\eta,N}(u)-u^{N+1}Ry_N.
\label{mem:identidad-finita}
\end{equation}
\end{proposition}
\begin{proof}
En grados $1\leq m\leq N$, el coeficiente del miembro izquierdo es
$y_m-Ry_{m-1}=\eta_{m-1}$, por \eqref{eq:eta}.
Sus coeficientes en grados cero y $N+1$ son $y_0$ y $-Ry_N$,
respectivamente. La comparación de coeficientes prueba la igualdad.
\end{proof}

El término terminal conserva la memoria transportada al cortar una historia.
La identidad permite comparar una evaluación finita con su prolongación
manteniendo explícito el estado que queda en la frontera.

\subsubsection{Correspondencia con el transporte del núcleo común}

La numeración de ventanas de la sección~\ref{subsec:memoria-resolvente}
y la empleada aquí describen la misma partición de la historia. La
correspondencia es
\[
 \underbrace{I_m=E_{m+1}}_{\text{índice del núcleo común}}
 \quad\longleftrightarrow\quad
 \underbrace{I_m=E_m}_{\text{índice de esta sección}},
 \qquad 0\leq m<12.
\]
El subíndice cambia el nombre de la ventana; sus nueve transiciones,
su orientación \(\sigma_m=\Sigma_{12}(m)\), su evento \(a_m\) y
su posición en el registro permanecen iguales. Para toda prolongación,
\(R=S^{-1}\), \(y_m=S^{-m}\mathbf z_m\) y la
actualización~\eqref{eq:eta} coinciden con
\eqref{mem:actualizacion}. La orientación y la secuencia de eventos
conservan la dependencia respecto de la historia, también después del
retorno del marco.

Esta identificación permite aplicar directamente la
proposición~\ref{mem:prop-serie}: la serie~\eqref{mem:serie}, la
cota~\eqref{mem:cota} y la convergencia uniforme de sus derivadas
corresponden a este mismo registro y al mismo reloj de ventanas. La
identidad finita~\eqref{mem:identidad-finita} determina además el término
terminal al truncar esa serie. Sus hipótesis mantienen la historia inicial
y el evento acotado \(|a_m|\leq9\).

Al reutilizar el registro, la proposición~\ref{mem:prop-schur} conserva
conjuntamente el operador, la fuente y el lector mediante
\eqref{mem:schur-fuente} y~\eqref{mem:schur-lector}. El ejemplo
\eqref{mem:ejemplo-ciclo} precisa el efecto de las excursiones omitidas
por una compresión. La composición afín~\eqref{mem:cociclo-tres}, los
bloques intermedios~\eqref{mem:schur-intermedio} y la identidad
\eqref{mem:schur-transitivo} garantizan la coherencia al concatenar
segmentos o eliminar varias capas. Sus demostraciones completas residen
en la sección~\ref{subsec:memoria-resolvente}, con los dominios, las
fuentes y las aplicaciones de lectura allí especificados.

El balance~\eqref{ii:mem:balance} conserva la correspondencia entre los
lectores orientados del registro firmado y la forma cuadrática común.
La distinción entre el vector acumulado \(\mathbf z_m\) y el bloque
decimal incidencial \(z_m\) sigue siendo la establecida en esta sección.
Las identidades de transporte se aplican a la memoria así tipada; no
determinan por sí solas coeficientes analíticos adicionales.
